\subsection{概论}

回顾 (参见 TG, IX, p.~37, 定义 9)，赋范代数是指 K 上的一个代数 A，它带有一个范数 \(x \mapsto \|x\|\)，使得：

\[
\| x y \| \leqslant \| x \| \| y \|\tag{1}
\]

对一切 \(x, y \in A\) 成立。若 A 完备，则称 A 是一个 Banach 代数。

设 A 是 K 上的一个完备可赋范代数。A 的拓扑可由一个满足 (1) 的范数定义 (参见 TG, IX, p.~38)。当 A 带有由这样一个范数所定义的赋范代数结构时，我们也说代数 A 是一个 Banach 代数。

我们还回顾下列事实 (参见 TG, IX, p.~38--39)：

\begin{enumerate}
\def\labelenumi{(\arabic{enumi})}
\item
  若 A 是一个非零赋范代数且有单位元 e，则 \(\|e\|\geqslant1\)。
\item
  设 A 是一个赋范 K-代数。A 的反代数，带上同一个范数，是一个赋范代数。A 的完备化，带上由 A 的范数连续延拓而得的范数，是一个 K-Banach 代数。A 的每个子代数，带上诱导范数，是一个赋范代数。若 I 是 A 的一个闭双边理想，则代数 A/I，带上由 \(\|\dot{x}\| = \inf_{x \in \dot{x}} \|x\|\)（对一切 \(\dot{x} \in A/I\)）所定义的范数，是一个赋范代数。
\item
  设 \((\mathrm{A}_{i})_{i\in\mathrm{I}}\) 是一族赋范代数，B 是诸代数 \(A_{i}\) 的积代数。由 B 的元素 \((x_{i})_{i\in\mathrm{I}}\)（满足 \(\|(\boldsymbol{x}_{i})\|=\sup_{i\in\mathrm{I}}\|\boldsymbol{x}_{i}\|<+\infty\)）所成的子代数是一个赋范代数，称为诸代数 \((\mathrm{A}_{i})_{i\in\mathrm{I}}\) 的积赋范代数。若对每个 i，\(A_{i}\) 都是 Banach 代数，则 B 是 Banach 代数。
\end{enumerate}

设 A 是一个赋范代数。设 \((y_{i})_{i\in\mathrm{I}}\) 是 A 的元素的一个族；包含诸元素 \(y_{i}\) 的最小闭子代数 B 称为由诸 \(y_{i}\) 生成的 A 的闭子代数；若 B = A，就说诸 \(y_{i}\) 拓扑地生成赋范代数 A，或者说族 \((y_{i})_{i\in\mathrm{I}}\) 是赋范代数 A 的一个拓扑生成系。

类似地，若 A 是一个赋范幺代数，则包含诸元素 \(y_{i}\) 的最小闭幺子代数称为由诸 \(y_{i}\) 生成的 A 的闭幺子代数。若它等于 A，就说诸 \(y_{i}\) 拓扑地生成赋范幺代数 A。

设 A 是一个赋范 K-代数。用 \(\tilde{\mathbf{A}}\) 表示由 A 添加单位元而得到的代数。在 \(\tilde{\mathbf{A}}\) 上，令 \(\| (\lambda ,x)\| = |\lambda | + \| x\|\)（对一切 \(\lambda \in \mathbf{K}\) 与每个 \(x\in \mathbf{A}\)），这样就定义了一个范数。我们有

\[
\begin{array}{r l} & {\| (\lambda , x) (\mu , y) \| = | \lambda \mu | + \| x y + \mu x + \lambda y \|} \\ & {\qquad \leqslant | \lambda | | \mu | + \| x \| \| y \| + | \mu | \| x \| + | \lambda | \| y \|} \\ & {\qquad = \| (\lambda , x) \| \| (\mu , y) \|.} \end{array}
\]

于是 \(\widetilde{\mathsf{A}}\) 成为一个赋范代数，称为由 A 添加单位元而得到的赋范幺代数。代数 A 与闭双边理想 \(\{0\} \times \mathbf{A}\)（\(\widetilde{\mathsf{A}}\) 的闭双边理想）视为同一。

定义 1. --- 设 A 是一个赋范代数。对 \(a \in \mathrm{A}\)，用 \(\gamma_{a}\) 与 \(\delta_{a}\) 表示从 A 到 A 的映射 \(x \mapsto ax\) 与 \(x \mapsto xa\)。映射 \(\gamma: a \mapsto \gamma_{a}\) 是 A 在 A 中的一个表示，称为 A 的左正则表示。映射 \(\delta: a \mapsto \delta_{a}\) 是 A 的反代数在 A 中的一个表示，称为 A 的右正则表示。

引理 1. --- 设 A 是一个赋范代数。A 的左正则表示与 A 的右正则表示都是连续的，且范数 \(\leqslant 1\)。

若代数 A 是幺的，则映射 \(x \mapsto \|\gamma_{x}\|\) 与 \(x \mapsto \|\delta_{x}\|\) 是 A 上的范数，它们定义 A 的拓扑。它们满足不等式 (1)。

若赋范幺代数 A 非零，也就是说，若 \(e \neq 0\)，则我们还有 \(\|\gamma_{e}\| = \|\delta_{e}\| = 1\)。

立即有 \(\|\gamma_{x}\|\leqslant\|x\|\) 与 \(\|\delta_{x}\|\leqslant\|x\|\)。

若 A 有单位元 \(e\)，则我们有 \(x = \gamma_x e = \delta_x e\)，因而：

\[
\| x \| \leqslant \| \boldsymbol {\gamma} _ {x} \| \cdot \| e \| \quad \| x \| \leqslant \| \boldsymbol {\delta} _ {x} \| \cdot \| e \|.\tag{2}
\]

这时，\(x \mapsto \|\gamma_{x}\|\) 与 \(x \mapsto \|\delta_{x}\|\) 是与 A 的范数等价的范数。它们满足 (1)，因为 \(\gamma\) 与 \(\delta\) 是表示。

最后一个断言由 \(\gamma_{e} = \delta_{e} = Id_{A}\) 这一事实得出。

